% ===========================================================================
%  1 — Introducao
% ===========================================================================
\chapter{Introdução}
\label{cap:introducao}

\section{Contexto e motivação}

Máquinas rotativas --- motores elétricos, bombas, ventiladores e redutores ---
concentram a maior parte dos ativos críticos de uma planta industrial. A falha
desses equipamentos raramente é instantânea: ela se anuncia por semanas em
sinais de vibração, temperatura e corrente, e a leitura correta desses sinais é
o que separa uma parada programada de \SI{2}{\hour} de uma parada não programada
de \SI{2}{\day}.

A literatura de manutenção organiza essa capacidade em uma escala de maturidade
crescente, e a distinção entre os degraus é operacional, não apenas
terminológica~\cite{abnt5462, jardine2006review, zonta2020predictive}:

\begin{description}
  \item[Corretiva] atua depois da falha. Custo máximo, previsibilidade nula.
  \item[Preventiva] atua em intervalos fixos, independentemente da condição
    real. Troca peças boas e ainda assim é surpreendida por falhas precoces.
  \item[Preditiva] monitora a condição e estima \emph{quando} a falha
    ocorrerá~\cite{iso17359, iso13381}. Responde \enquote{este rolamento vai
    falhar}, mas deixa o \enquote{e agora?} para o julgamento do técnico.
  \item[Prescritiva] acrescenta a resposta que falta: \emph{o que fazer},
    \emph{em que ordem} e \emph{com base em qual procedimento
    documentado}~\cite{ansari2019prima}.
\end{description}

O salto da manutenção preditiva para a prescritiva é o que este projeto
implementa. Um diagnóstico que apenas classifica a falha transfere ao operador
o trabalho de localizar, no acervo de manuais da empresa, qual procedimento se
aplica --- e é exatamente aí que o conhecimento organizacional se perde. A
solução documentada aqui fecha esse ciclo: identifica o tipo de defeito,
quantifica quantas vezes ele já ocorreu no histórico, e devolve as instruções de
correção extraídas dos documentos técnicos da própria planta, com citação de
documento, seção e página.

\subsection{O problema proposto}

O enunciado do processo seletivo estabelece o seguinte problema: dado um novo
evento de sensores de vibração em formato JSON, o sistema deve

\begin{enumerate}
  \item identificar o \textbf{tipo de defeito} presente no evento;
  \item quantificar as \textbf{ocorrências históricas semelhantes} ---
    quantidade, distribuição temporal e frequência;
  \item gerar as \textbf{instruções de correção} consultando a base documental
    da empresa;
  \item restringir-se \textbf{apenas a falhas documentadas}: quando o defeito
    identificado não possui documento orientativo, informar explicitamente essa
    condição e sugerir o registro de um novo documento.
\end{enumerate}

O quarto requisito é o mais exigente do ponto de vista de engenharia. Ele
determina que o sistema saiba, com certeza verificável, o que ele \emph{não}
sabe. Modelos de linguagem são notoriamente ruins nisso: quando não encontram
a informação, tendem a produzir uma resposta plausível em vez de
admitir a lacuna~\cite{ji2023hallucination}. A resposta arquitetural adotada foi
retirar essa decisão do modelo: a cobertura documental é verificada por um
\textit{guardrail} determinístico antes de qualquer chamada ao modelo de
linguagem, conforme detalhado no Capítulo~\ref{cap:rag}.

\section{Objetivo do documento}

Este documento é a especificação técnica completa da solução. Ele foi escrito
para ser suficiente por si só: um engenheiro que nunca viu o repositório deve
conseguir, apenas com este texto, compreender o problema, avaliar as decisões de
projeto, reproduzir o ambiente, operar o sistema e modificar o código com
segurança.

Especificamente, o documento cumpre os seguintes objetivos:

\begin{itemize}
  \item \textbf{Delimitar o escopo} --- o que a solução faz, o que
    deliberadamente não faz, sob quais premissas e restrições
    (Capítulo~\ref{cap:escopo}).
  \item \textbf{Descrever a arquitetura} segundo a norma
    ISO/IEC/IEEE 42010~\cite{iso42010}, usando o modelo
    C4~\cite{brown2026c4} nos quatro níveis e registrando as decisões em
    formato ADR~\cite{nygard2011adr} (Capítulo~\ref{cap:arquitetura}).
  \item \textbf{Especificar o comportamento} por meio de casos de uso UML
    detalhados (Capítulo~\ref{cap:casosdeuso}) e de diagramas de atividade de
    todos os processos do sistema (Capítulo~\ref{cap:atividades}).
  \item \textbf{Fundamentar as escolhas de aprendizado de máquina} com
    referência à literatura, incluindo a justificativa do algoritmo, o protocolo
    de validação adotado e o diagnóstico honesto de suas limitações
    (Capítulo~\ref{cap:ml}).
  \item \textbf{Documentar o código} módulo a módulo, com propósito,
    interface pública, invariantes e pontos de extensão
    (Capítulo~\ref{cap:codigo}).
\end{itemize}

\begin{nota}
Este documento descreve o sistema \emph{como ele foi construído}, não como se
gostaria que ele fosse. Onde os resultados são desfavoráveis --- notadamente na
generalização do classificador, Seção~\ref{sec:vazamento} --- eles são
reportados com o número real, a causa diagnosticada e o encaminhamento
proposto. Um documento técnico que esconde o resultado ruim perde a função de
apoiar decisão.
\end{nota}

\section{Público-alvo e formas de leitura}

\begin{tabularx}{\linewidth}{@{}L{34mm}L{40mm}X@{}}
\toprule
\headrow \thd{Perfil} & \thd{Interesse principal} & \thd{Roteiro sugerido}\\
\midrule
Avaliador técnico & Aderência ao enunciado e qualidade de engenharia &
  Capítulos~\ref{cap:escopo}, \ref{cap:arquitetura}, \ref{cap:ml}
  e~\ref{cap:rag}\\
\zebra Arquiteto de \textit{software} & Estrutura, fronteiras e
  \textit{trade-offs} &
  Capítulos~\ref{cap:arquitetura}, \ref{cap:modelodados}
  e~\ref{cap:infraestrutura}\\
Cientista de dados & Dados, \textit{features}, modelo e validação &
  Capítulos~\ref{cap:pipeline}, \ref{cap:ml} e Apêndice~\ref{ap:features}\\
\zebra Desenvolvedor & Como o código está organizado e como estendê-lo &
  Capítulos~\ref{cap:api}, \ref{cap:codigo} e Apêndice~\ref{ap:config}\\
Engenheiro de manutenção & O que o sistema responde e com que confiança &
  Capítulos~\ref{cap:casosdeuso}, \ref{cap:interface} e
  Seção~\ref{sec:confianca}\\
\zebra Equipe de operação & Implantar, monitorar e diagnosticar falhas &
  Capítulos~\ref{cap:infraestrutura}, \ref{cap:qualidade} e
  Apêndice~\ref{ap:operacao}\\
\bottomrule
\end{tabularx}

\section{Convenções adotadas}

\subsection{Tipografia e identificadores}

\begin{itemize}
  \item \arq{caminho/arquivo.py} indica arquivo ou diretório do repositório.
  \item \cd{identificador} indica nome de função, classe, variável, tabela,
    coluna, tópico MQTT ou variável de ambiente.
  \item Termos técnicos sem tradução consolidada aparecem em \textit{itálico}
    na primeira ocorrência.
  \item Identificadores rastreáveis seguem prefixos fixos:
    \req{RF-nn} (requisito funcional), \req{RNF-nn} (requisito não funcional),
    \ucr{UC-nn} (caso de uso), \adrr{ADR-nn} (decisão arquitetural),
    \req{AT-nn} (diagrama de atividade) e \req{RI-nn} (risco).
\end{itemize}

\subsection{Diagramas}

Todos os diagramas são vetoriais, desenhados nativamente em \TeX{} com a
biblioteca TikZ, o que garante tipografia uniforme e nitidez em qualquer nível
de ampliação. O código-fonte \textbf{Mermaid} equivalente de cada diagrama ---
utilizável diretamente em \arq{README.md}, GitHub, GitLab ou ferramentas de
documentação --- está reunido no Apêndice~\ref{ap:mermaid}, indexado pelo mesmo
número de figura.

As notações empregadas são:

\begin{tabularx}{\linewidth}{@{}L{42mm}X@{}}
\toprule
\headrow \thd{Notação} & \thd{Onde é usada}\\
\midrule
Modelo C4~\cite{brown2026c4} & Contexto, contêineres, componentes e código
  (Capítulo~\ref{cap:arquitetura})\\
\zebra UML --- casos de uso~\cite{omg2017uml} & Capítulo~\ref{cap:casosdeuso}\\
UML --- atividades~\cite{omg2017uml} & Capítulo~\ref{cap:atividades} e
  processos ao longo do texto\\
\zebra UML --- sequência~\cite{omg2017uml} & Fluxos de interação
  (Capítulos~\ref{cap:arquitetura}, \ref{cap:rag} e~\ref{cap:mqtt})\\
UML --- classes~\cite{omg2017uml} & Nível 4 do C4
  (Seção~\ref{sec:c4nivel4})\\
\zebra Entidade-Relacionamento & Modelo de dados
  (Capítulo~\ref{cap:modelodados})\\
Diagrama de implantação & Infraestrutura local e em nuvem
  (Capítulo~\ref{cap:infraestrutura})\\
\bottomrule
\end{tabularx}

\subsection{Números e unidades}

Todas as grandezas físicas são expressas em unidades do SI, com vírgula como
separador decimal e ponto como separador de milhar, conforme a prática
brasileira. As colunas redundantes do conjunto de dados original em
\si{\inch\per\second} e \si{\degF} foram descartadas do vetor de
\textit{features} por colinearidade exata com suas contrapartes métricas
(Seção~\ref{sec:features}).

Métricas do modelo são reproduzidas exatamente como registradas em
\arq{artifacts/metrics.json}, gerado pela execução de treino de
22 de agosto de 2026 sobre \num{166796} eventos.

\section{Estrutura do documento}

\begin{tabularx}{\linewidth}{@{}L{13mm}L{46mm}X@{}}
\toprule
\headrow \thd{Cap.} & \thd{Título} & \thd{Conteúdo}\\
\midrule
\ref{cap:introducao} & Introdução & Contexto, objetivo, público e convenções.\\
\zebra \ref{cap:definicoes} & Definições & Glossário de domínio e de
  tecnologia.\\
\ref{cap:escopo} & Escopo & Problema, fronteiras, premissas, requisitos e
  rastreabilidade.\\
\zebra \ref{cap:casosdeuso} & Casos de uso & Atores, diagrama UML e
  especificação de doze casos de uso.\\
\ref{cap:arquitetura} & Arquitetura & Modelo C4 completo, ADRs e atributos de
  qualidade.\\
\zebra \ref{cap:modelodados} & Modelo de dados & Diagrama ER, dicionário de
  dados e esquemas.\\
\ref{cap:pipeline} & Pipeline de dados & ETL, canonicalização de rótulos e
  auditoria.\\
\zebra \ref{cap:ml} & Aprendizado de máquina & \textit{Features}, escolha do
  algoritmo, treino, validação e limitações.\\
\ref{cap:rag} & RAG e linguagem & Ingestão documental, \textit{guardrail},
  \textit{prompting} e agente.\\
\zebra \ref{cap:mqtt} & Integração industrial & Tópicos MQTT, garantias de
  entrega, DLQ e alertas.\\
\ref{cap:interface} & Interface & Navegação, painéis analíticos e decisões de
  visualização.\\
\zebra \ref{cap:api} & API & Especificação de todos os \textit{endpoints}
  REST.\\
\ref{cap:atividades} & Atividades & Diagramas de atividade de todos os
  processos.\\
\zebra \ref{cap:infraestrutura} & Infraestrutura & Ambiente local, nuvem, IaC,
  segredos e custos.\\
\ref{cap:qualidade} & Qualidade & Estratégia de testes, suíte, CI/CD e padrões
  de código.\\
\zebra \ref{cap:codigo} & Documentação de código & Referência de todos os
  módulos.\\
\ref{cap:riscos} & Riscos e evolução & Matriz de riscos e \textit{roadmap}
  técnico.\\
\zebra \ref{cap:conclusao} & Conclusão & Síntese e aderência ao enunciado.\\
\bottomrule
\end{tabularx}
